\documentclass[a4paper]{article}
\usepackage{amsfonts}
\usepackage{amsmath}
\usepackage{amssymb}
\usepackage{graphicx}

\begin{document}
  
\noindent \large Auteur: Siebe Mees \\
\noindent \large Studierichting: Informatica\\
\noindent \large Oefeningenreeks 4A nr. 43.6\\

\medskip

\normalsize

$ \displaystyle\int \dfrac{9}{4x^3-3x+1}dx $\\

Probeer $\dfrac{1}{1,2,4}$ , nulpunt bij $x = \dfrac{1}{2} $\\

$ \displaystyle\int \dfrac{9}{(2x-1)(2x^2+x-1)}dx $\\

$ \Delta =1^2-4\cdot 2\cdot (-1)=9 \qquad x_{1,2}=\dfrac{-1\pm \sqrt{9}}{2\cdot 2} \sim x_1=\dfrac{1}{2}, \ x_2=-1 $\\

$ \displaystyle\int \dfrac{9}{(2x-1)^2(x+1)}dx=\dfrac{A}{2x-1}+\dfrac{B}{(2x-1)^2}+\dfrac{C}{x+1} $\\

Dan is \\

$ B=\dfrac{9}{x+1}\bigg|_{x=\tfrac{1}{2}}=6 $\\

$ C=\dfrac{9}{(2x-1)^2}\bigg|_{x=-1}=1 $\\

$ \left \{
\begin{array}{l}
\dfrac{9}{(2x-1)^2(x+1)} - \dfrac{B}{(2x-1)^2} = \dfrac{9}{(2x-1)^2(x+1)} - \dfrac{6}{(2x-1)^2} \\

= \dfrac{9 - 6(x+1)}{(2x-1)^2(x+1)} = \dfrac{-6x + 3}{(2x-1)^2(x+1)} \\

= \dfrac{-3(2x-1)}{(2x-1)^2(x+1)} = \dfrac{-3}{(2x-1)(x+1)} \\

\end{array}
\right . $\\

$ A = \dfrac{-3}{x+1} \bigg|_{x = \tfrac{1}{2}} = -2 $\\

$ \Rightarrow \displaystyle\int \left( \dfrac{-2}{2x-1}+\dfrac{6}{(2x-1)^2}+\dfrac{1}{x+1} \right)dx $\\

$ =\displaystyle\int \dfrac{-2}{2x-1}dx+\displaystyle\int \dfrac{6}{(2x-1)^2}dx+\displaystyle\int \dfrac{1}{x+1}dx $\\

$ t = 2x-1 \left(dt = 2dx, dx = \dfrac{dt}{2} \right)$\\

$ =\displaystyle\int \dfrac{-2}{2x-1}dx+6\displaystyle\int t^{-2}\dfrac{dt}{2}+\displaystyle\int \dfrac{1}{x+1}dx $\\

$ =\displaystyle\int \dfrac{-2}{2x-1}dx+3\displaystyle\int t^{-2}dt+\displaystyle\int \dfrac{1}{x+1}dx $\\

$ =-\ln|2x-1|-\dfrac{3}{2x-1}+\ln|x+1|+C $\\

\begin{thebibliography}{999}
%\bibitem{a} Referentie
\end{thebibliography}
\end{document} 
